\section{Proof of Cauchy-Schwarz Inequality}

\textbf{(a)}

We need to check that its minimum value is also nonnegative.
The quadratic is minimized at
\[
    \lambda^\star = -\frac{b}{c}
\]
Substituting gives
\[
    a + 2b\left(-\frac{b}{c}\right)
        + c\left(-\frac{b}{c}\right)^2
    =
    a - \frac{b^2}{c} \geq 0 \quad \implies b^2 \leq ac \quad \implies |b| \leq \sqrt{ac}
\]

\textbf{(b)}

%For any \(v,w \in \mathbf{R}^n\) and any \(\lambda \in \mathbf{R}\), the
%vector \(v+\lambda w\) is in \(\mathbf{R}^n\).
The inner product of any vector with itself is the sum of the squares of its entries, so it is
always nonnegative:
\[
    (v+\lambda w)^T(v+\lambda w)
    =
    \|v+\lambda w\|^2
    \geq 0
\]

\textbf{(c)}
\begin{align*}
    (v+\lambda w)^T(v+\lambda w)
    &=
    v^T v + 2(v^T w)\lambda + (w^T w)\lambda^2
\end{align*}

This has the same form as part (a), with
\[
    a = v^T v,
    \qquad
    b = v^T w,
    \qquad
    c = w^T w
\]

So
\[
    |v^T w|
    \leq
    \sqrt{(v^T v)(w^T w)}
    =
    \sqrt{v^T v}\sqrt{w^T w}
\]

\textbf{(d)}

Equality holds iff \(v\) and \(w\) are linearly dependent

For example, with $v = [1 \quad 0]^\top$ and $v = [2 \quad 0]^\top$:
\begin{align*}
    |v^T w|
    &=
    \sqrt{v^T v}\sqrt{w^T w} \\
    2 &= 1 \cdot 2 = 2
\end{align*}

% To see this, if \(w \neq 0\), equality in part (a) means the quadratic
% \[
%     \|v+\lambda w\|^2
% \]
% has minimum value zero. Hence there is some \(\lambda\) such that
% \[
%     v+\lambda w = 0,
% \]
% so \(v\) is a scalar multiple of \(w\). Conversely, if \(v=\alpha w\),
% then
% \[
%     |v^T w|
%     =
%     |\alpha|\, w^T w
%     =
%     \sqrt{\alpha^2 w^T w}\sqrt{w^T w}
%     =
%     \sqrt{v^T v}\sqrt{w^T w},
% \]
% so equality holds.
